\documentclass[12pt,a4paper]{article}

% ==================== GÓI CƠ BẢN ====================
\usepackage[utf8]{inputenc}
\usepackage[T5]{fontenc}
\usepackage[vietnamese]{babel}
\usepackage{geometry}
\usepackage{parskip}
\usepackage{enumitem}
\usepackage{longtable}
\usepackage{array}
\usepackage{xcolor}
\usepackage{hyperref}
\usepackage{fancyhdr}
\usepackage{lastpage}

% ==================== ĐỊNH DẠNG TRANG ====================
\geometry{
    top=2cm,
    bottom=2cm,
    left=2.5cm,
    right=2cm
}

\setlength{\parindent}{0pt}
\setlength{\parskip}{0.7em}
\setlist[itemize]{leftmargin=1.8em, itemsep=0.35em}
\setlist[enumerate]{leftmargin=1.8em, itemsep=0.35em}

\definecolor{ailmsblue}{RGB}{25,84,166}
\definecolor{ailmsgray}{RGB}{90,90,90}
\definecolor{ailmsred}{RGB}{180,40,40}

\hypersetup{
    colorlinks=true,
    linkcolor=ailmsblue,
    urlcolor=ailmsblue,
    citecolor=ailmsblue,
    pdftitle={Chính sách Nhân sự và Giảng viên AILMS},
    pdfauthor={Phòng Nhân sự AILMS}
}

% ==================== HEADER / FOOTER ====================
\pagestyle{fancy}
\fancyhf{}
\fancyhead[L]{\textcolor{ailmsgray}{AILMS}}
\fancyhead[R]{\textcolor{ailmsgray}{Chính sách Nhân sự \& Giảng viên}}
\fancyfoot[C]{Trang \thepage/\pageref{LastPage}}
\renewcommand{\headrulewidth}{0.4pt}
\renewcommand{\footrulewidth}{0.4pt}

% ==================== THÔNG TIN TÀI LIỆU ====================
\title{
    \textbf{\LARGE CHÍNH SÁCH NHÂN SỰ\\
    VÀ QUẢN LÝ GIẢNG VIÊN AILMS}
}
\author{Phòng Nhân sự AILMS}
\date{Phiên bản 1.0 -- Ngày hiệu lực: \today}

\begin{document}

\maketitle

\begin{center}
    \textit{Tài liệu này thiết lập các nguyên tắc tuyển dụng, xác minh,
    phân quyền, quản lý công việc, đánh giá và bảo vệ quyền lợi
    của nhân sự, giảng viên và trợ giảng trên hệ thống AILMS.}
\end{center}

\vspace{0.8cm}

\begin{longtable}{|>{\bfseries}p{4cm}|p{9.5cm}|}
    \hline
    Tên tài liệu &
    Chính sách Nhân sự và Quản lý Giảng viên AILMS \\
    \hline
    Phiên bản & 1.0 \\
    \hline
    Đơn vị quản lý & Phòng Nhân sự AILMS \\
    \hline
    Đối tượng áp dụng &
    Giảng viên, trợ giảng, nhân viên toàn thời gian, nhân viên bán thời gian,
    cộng tác viên, nhân sự quản lý và quản trị viên có liên quan \\
    \hline
    Căn cứ thực hiện &
    Hợp đồng đã ký, quy chế nội bộ và pháp luật Việt Nam hiện hành \\
    \hline
    Chu kỳ rà soát &
    Tối thiểu 12 tháng một lần hoặc khi có thay đổi về pháp luật,
    cơ cấu tổ chức hay chức năng hệ thống \\
    \hline
\end{longtable}

\tableofcontents
\newpage

% =========================================================
\section{Mục đích}

Chính sách này được ban hành nhằm:

\begin{itemize}
    \item Thiết lập quy trình tiếp nhận và xác minh hồ sơ nhân sự.
    \item Quy định nguyên tắc cấp tài khoản, vai trò và quyền truy cập.
    \item Xác định trách nhiệm chuyên môn của giảng viên và trợ giảng.
    \item Bảo đảm minh bạch trong phân công, chấm công, tiền lương và đánh giá.
    \item Bảo vệ dữ liệu cá nhân, hồ sơ lao động và thông tin thu nhập.
    \item Phòng ngừa xung đột lợi ích, quấy rối và hành vi không phù hợp.
    \item Thiết lập cơ chế giải trình, khiếu nại và xử lý vi phạm.
\end{itemize}

Chính sách này không thay thế hợp đồng lao động, hợp đồng cộng tác,
thỏa thuận giảng dạy hoặc quy định bắt buộc của pháp luật.

Nếu có khác biệt, văn bản có giá trị pháp lý cao hơn và thỏa thuận hợp lệ
được ưu tiên áp dụng.

% =========================================================
\section{Giải thích thuật ngữ}

\begin{itemize}
    \item \textbf{Nhân sự:} Cá nhân làm việc hoặc cung cấp dịch vụ cho AILMS
    theo hợp đồng hoặc thỏa thuận hợp lệ.

    \item \textbf{Giảng viên -- TEACHER:} Người được giao nhiệm vụ
    xây dựng khóa học, giảng dạy, đánh giá và hỗ trợ học viên.

    \item \textbf{Trợ giảng -- TA:} Người hỗ trợ giảng viên trong phạm vi
    chuyên môn và quyền hạn được giao.

    \item \textbf{Nhân sự -- HR:} Người thực hiện công tác tuyển dụng,
    hồ sơ, hợp đồng, chấm công, hỗ trợ và quản trị nhân sự.

    \item \textbf{Quản trị viên -- ADMIN:} Người quản lý cấu hình,
    dữ liệu hoặc hoạt động hệ thống theo quyền được cấp.

    \item \textbf{Employee Code:} Mã định danh nội bộ của nhân sự.
    Mã này không thay thế giấy tờ tùy thân hoặc mã số do cơ quan nhà nước cấp.

    \item \textbf{Teaching Session:} Buổi giảng dạy hoặc tư vấn
    đã được lập lịch và ghi nhận trên hệ thống.

    \item \textbf{Teaching Rate:} Đơn giá giảng dạy được thỏa thuận
    cho một giờ, một buổi hoặc một khối lượng công việc.

    \item \textbf{Audit Log:} Nhật ký ghi nhận các thao tác quan trọng
    phục vụ kiểm tra và truy vết.
\end{itemize}

% =========================================================
\section{Nguyên tắc quản lý nhân sự}

AILMS áp dụng các nguyên tắc:

\begin{itemize}
    \item Tuyển dụng và phân công dựa trên năng lực, yêu cầu công việc
    và tiêu chí đã công bố.
    \item Không phân biệt đối xử trái pháp luật.
    \item Tôn trọng danh dự, nhân phẩm và quyền riêng tư của người lao động.
    \item Chỉ thu thập dữ liệu nhân sự cần thiết và có mục đích rõ ràng.
    \item Phân quyền theo nguyên tắc tối thiểu cần thiết.
    \item Đánh giá dựa trên nhiều nguồn thông tin, không chỉ dựa vào thuật toán.
    \item Tiền lương, quyền lợi và trách nhiệm phải được xác định bằng
    hợp đồng hoặc quy chế hợp lệ.
\end{itemize}

% =========================================================
\section{Tuyển dụng và tiếp nhận hồ sơ}

\subsection{Hồ sơ cần cung cấp}

Tùy theo vị trí, ứng viên có thể được yêu cầu cung cấp:

\begin{itemize}
    \item Thông tin định danh và liên hệ.
    \item Sơ yếu lý lịch hoặc hồ sơ năng lực.
    \item Bằng cấp, chứng chỉ và minh chứng chuyên môn.
    \item Kinh nghiệm giảng dạy hoặc làm việc.
    \item Danh mục công trình, khóa học hoặc tài liệu đã thực hiện.
    \item Thông tin phục vụ ký kết hợp đồng và thanh toán.
    \item Thông tin khác cần thiết cho vị trí tuyển dụng.
\end{itemize}

AILMS không yêu cầu ứng viên cung cấp dữ liệu không liên quan đến công việc
nếu không có căn cứ phù hợp.

\subsection{Xác minh hồ sơ}

Hồ sơ có thể được xác minh thông qua:

\begin{itemize}
    \item Đối chiếu giấy tờ và thông tin do ứng viên cung cấp.
    \item Xác minh bằng cấp hoặc chứng chỉ với đơn vị cấp khi cần thiết.
    \item Phỏng vấn chuyên môn.
    \item Giảng thử hoặc thực hiện bài đánh giá năng lực.
    \item Kiểm tra người tham chiếu sau khi có căn cứ phù hợp.
\end{itemize}

Ứng viên phải được thông báo khi AILMS cần thực hiện hoạt động xác minh
có liên quan đến bên thứ ba hoặc dữ liệu cá nhân nhạy cảm.

\subsection{Thông tin không chính xác}

Nếu phát hiện hồ sơ có thông tin sai lệch, AILMS phải tạo điều kiện
để ứng viên hoặc nhân sự giải trình trước khi đưa ra quyết định,
trừ trường hợp cần áp dụng biện pháp tạm thời để bảo vệ hệ thống.

Việc làm giả bằng cấp, giấy tờ hoặc kinh nghiệm có thể dẫn đến:

\begin{itemize}
    \item Từ chối tiếp nhận.
    \item Tạm dừng phân công.
    \item Thu hồi quyền truy cập liên quan.
    \item Xử lý theo hợp đồng và quy định pháp luật.
\end{itemize}

% =========================================================
\section{Ký kết và quản lý hợp đồng}

Trước khi bắt đầu công việc chính thức, các bên phải xác lập quan hệ
bằng loại hợp đồng hoặc thỏa thuận phù hợp.

Hợp đồng nên xác định tối thiểu:

\begin{itemize}
    \item Danh tính và thông tin của các bên.
    \item Vị trí, nhiệm vụ và phạm vi công việc.
    \item Thời hạn và địa điểm làm việc.
    \item Thời giờ làm việc hoặc nguyên tắc phân công.
    \item Tiền lương, đơn giá, phụ cấp và phương thức thanh toán.
    \item Quyền sở hữu trí tuệ.
    \item Nghĩa vụ bảo mật.
    \item Điều kiện sửa đổi, tạm hoãn và chấm dứt.
    \item Trách nhiệm đối với thiết bị, tài khoản và dữ liệu.
\end{itemize}

Trạng thái hợp đồng trên hệ thống phải phản ánh hồ sơ nghiệp vụ,
nhưng không tự thay thế văn bản hợp đồng đã ký.

Mọi thay đổi quan trọng về chức danh, mức lương, phạm vi công việc
hoặc thời hạn phải được ghi nhận bằng hình thức hợp lệ.

% =========================================================
\section{Phân loại nhân sự}

\subsection{Nhân sự toàn thời gian}

Nhân sự toàn thời gian thực hiện công việc theo thời giờ,
nhiệm vụ và chế độ được quy định trong hợp đồng.

Quyền lợi có thể bao gồm:

\begin{itemize}
    \item Tiền lương cơ bản.
    \item Phụ cấp và thưởng theo chính sách.
    \item Nghỉ phép và các chế độ nghỉ khác.
    \item Bảo hiểm và quyền lợi bắt buộc nếu thuộc đối tượng áp dụng.
    \item Đào tạo và đánh giá định kỳ.
\end{itemize}

\subsection{Nhân sự bán thời gian}

Nhân sự bán thời gian làm việc theo khung thời gian hoặc khối lượng
đã thỏa thuận và được hưởng quyền lợi theo hợp đồng cùng pháp luật áp dụng.

Việc làm bán thời gian không phải là căn cứ để đối xử bất lợi trái pháp luật
so với người làm toàn thời gian trong các quyền lợi có cùng điều kiện áp dụng.

\subsection{Giảng viên và trợ giảng theo buổi}

Giảng viên hoặc trợ giảng có thể được thanh toán dựa trên:

\begin{itemize}
    \item Số buổi giảng dạy được xác nhận.
    \item Số giờ giảng thực tế.
    \item Khối lượng chấm bài hoặc hỗ trợ.
    \item Mốc hoàn thành nội dung.
    \item Đơn giá và điều kiện trong hợp đồng.
\end{itemize}

Việc gọi một cá nhân là cộng tác viên hoặc giảng viên theo buổi
không tự quyết định bản chất pháp lý của quan hệ giữa các bên.

% =========================================================
\section{Tạo tài khoản và cấp mã nhân sự}

Sau khi hoàn tất thủ tục cần thiết, hệ thống có thể:

\begin{enumerate}
    \item Tạo hồ sơ nhân sự.
    \item Sinh Employee Code duy nhất.
    \item Liên kết hồ sơ nhân sự với tài khoản người dùng.
    \item Cấp vai trò và quyền phù hợp.
    \item Gửi hướng dẫn kích hoạt và bảo mật tài khoản.
\end{enumerate}

Employee Code phải:

\begin{itemize}
    \item Không trùng lặp.
    \item Không tái sử dụng tùy tiện cho người khác.
    \item Được sử dụng nhất quán trong hồ sơ nội bộ.
    \item Không chứa dữ liệu cá nhân nhạy cảm nếu không cần thiết.
\end{itemize}

% =========================================================
\section{Phân quyền và quản lý tài khoản}

\subsection{Nguyên tắc phân quyền}

Vai trò chỉ được cấp khi:

\begin{itemize}
    \item Có nhu cầu công việc hợp lệ.
    \item Đã hoàn tất phê duyệt theo quy trình.
    \item Hợp đồng hoặc phân công còn hiệu lực.
    \item Người dùng đã đáp ứng yêu cầu bảo mật.
\end{itemize}

Một tài khoản có thể có nhiều vai trò, nhưng mỗi quyền phải có căn cứ rõ ràng.

\subsection{Phạm vi vai trò}

\begin{itemize}
    \item \textbf{TEACHER:} Quản lý khóa học được giao, lớp học,
    lịch dạy, bài tập và đánh giá học viên.

    \item \textbf{TA:} Hỗ trợ lớp học, chấm bài, điểm danh và thảo luận
    trong phạm vi được giảng viên hoặc quản trị viên giao.

    \item \textbf{HR:} Quản lý hồ sơ, hợp đồng, chấm công, nghỉ phép
    và quy trình nhân sự theo thẩm quyền.

    \item \textbf{ADMIN:} Quản trị hệ thống theo phạm vi kỹ thuật
    hoặc nghiệp vụ được cấp.
\end{itemize}

Vai trò \texttt{ADMIN} không mặc nhiên cho phép đọc mọi dữ liệu nhân sự.
Dữ liệu tiền lương, hợp đồng và hồ sơ riêng tư phải có quyền chuyên biệt.

\subsection{Rà soát và thu hồi quyền}

Quyền truy cập phải được rà soát khi:

\begin{itemize}
    \item Nhân sự thay đổi vị trí.
    \item Hợp đồng hết hạn hoặc bị tạm hoãn.
    \item Nhân sự nghỉ việc.
    \item Phát hiện rủi ro bảo mật.
    \item Quyền không còn cần thiết.
\end{itemize}

Khi quan hệ công việc kết thúc, quyền nội bộ phải được thu hồi kịp thời,
đồng thời bảo toàn dữ liệu và nhật ký cần thiết.

% =========================================================
\section{Tiêu chuẩn nghề nghiệp của giảng viên}

Giảng viên phải:

\begin{itemize}
    \item Có chuyên môn phù hợp với nội dung giảng dạy.
    \item Cung cấp thông tin năng lực trung thực.
    \item Chuẩn bị nội dung theo mục tiêu khóa học.
    \item Thực hiện lịch giảng đã xác nhận.
    \item Đánh giá học viên công bằng và nhất quán.
    \item Tôn trọng sự khác biệt và quyền riêng tư.
    \item Không lợi dụng vị trí để trục lợi hoặc quấy rối.
    \item Không tiết lộ dữ liệu học viên trái phép.
    \item Cập nhật kiến thức và phương pháp giảng dạy khi cần thiết.
\end{itemize}

% =========================================================
\section{Tiêu chuẩn của trợ giảng}

Trợ giảng có trách nhiệm:

\begin{itemize}
    \item Hỗ trợ giảng viên theo phân công.
    \item Theo dõi câu hỏi và hoạt động lớp.
    \item Hỗ trợ điểm danh và xử lý bài nộp.
    \item Chấm bài theo rubric hoặc hướng dẫn đã được phê duyệt.
    \item Chuyển vấn đề vượt thẩm quyền cho giảng viên.
    \item Không tự ý thay đổi điểm hoặc điều kiện khóa học.
\end{itemize}

Trợ giảng không được thực hiện chức năng quản trị hoặc quyết định chuyên môn
vượt quá phạm vi được giao.

% =========================================================
\section{Xây dựng và phê duyệt khóa học}

\subsection{Tạo khóa học}

Giảng viên được phép tạo bản nháp khóa học nếu có quyền phù hợp.

Khóa học cần có:

\begin{itemize}
    \item Tên và mô tả rõ ràng.
    \item Mục tiêu học tập.
    \item Đối tượng và yêu cầu đầu vào.
    \item Cấu trúc nội dung.
    \item Tài liệu, bài tập và phương thức đánh giá.
    \item Thông tin bản quyền và nguồn tham khảo.
    \item Điều kiện hoàn thành.
\end{itemize}

\subsection{Quy trình phê duyệt}

Khóa học có thể trải qua các trạng thái:

\begin{enumerate}
    \item Bản nháp.
    \item Gửi duyệt.
    \item Đang xem xét.
    \item Yêu cầu chỉnh sửa.
    \item Được phê duyệt.
    \item Xuất bản.
    \item Tạm ẩn hoặc lưu trữ.
\end{enumerate}

Chỉ sử dụng trạng thái \texttt{PENDING\_APPROVAL} nếu enum này
thực sự tồn tại trong mã nguồn và cơ sở dữ liệu.

Việc phê duyệt cần xem xét:

\begin{itemize}
    \item Độ chính xác chuyên môn.
    \item Chất lượng và cấu trúc nội dung.
    \item Tính phù hợp với đối tượng học viên.
    \item Bản quyền và quyền sử dụng tài liệu.
    \item Nội dung không phù hợp hoặc có nguy cơ gây hại.
    \item Điều kiện kỹ thuật và khả năng vận hành.
\end{itemize}

Người duyệt phải ghi nhận lý do khi từ chối hoặc yêu cầu chỉnh sửa.

% =========================================================
\section{Tổ chức lớp học và lịch giảng}

Giảng viên có trách nhiệm:

\begin{itemize}
    \item Xác nhận lịch trước khi công bố.
    \item Theo dõi các buổi được phân công.
    \item Thông báo sớm nếu cần đổi hoặc hủy lịch.
    \item Ghi nhận nội dung đã giảng và tình trạng buổi học.
    \item Phối hợp bố trí lịch học bù khi lỗi thuộc phía giảng viên.
\end{itemize}

Đối với lớp một-một, sau khi trao đổi với học viên,
giảng viên hoặc trợ giảng có quyền phù hợp có thể tạo lớp và lịch học.

HR chỉ cần can thiệp khi:

\begin{itemize}
    \item Không tìm được giảng viên phù hợp.
    \item Có xung đột phân công.
    \item Giảng viên không thể tiếp tục.
    \item Có vấn đề về hợp đồng hoặc khiếu nại.
\end{itemize}

% =========================================================
\section{Chấm công và ghi nhận thời gian làm việc}

\subsection{Nguồn dữ liệu chấm công}

Hệ thống có thể ghi nhận:

\begin{itemize}
    \item Giờ vào và giờ ra.
    \item Lịch làm việc.
    \item Lịch giảng và trạng thái buổi học.
    \item Đơn nghỉ phép.
    \item Xác nhận của người quản lý.
    \item Điều chỉnh đã được phê duyệt.
\end{itemize}

\subsection{Trạng thái chấm công}

Các trạng thái có thể bao gồm:

\begin{itemize}
    \item Có mặt.
    \item Đi muộn.
    \item Về sớm.
    \item Vắng có phép.
    \item Vắng không phép.
    \item Nghỉ lễ.
    \item Nghỉ không hưởng lương.
    \item Chờ xác nhận.
\end{itemize}

\subsection{Nguyên tắc xử lý}

Dữ liệu tự động chỉ là nguồn hỗ trợ đối soát.

Hệ thống không được tự động khấu trừ tiền lương hoặc áp dụng kỷ luật
chỉ dựa trên một bản ghi đi muộn hay vắng mặt chưa được xác minh.

Trước khi khóa bảng công, nhân sự phải có cơ hội:

\begin{itemize}
    \item Xem dữ liệu liên quan.
    \item Gửi yêu cầu điều chỉnh.
    \item Cung cấp lý do hoặc bằng chứng.
    \item Nhận kết quả xử lý.
\end{itemize}

% =========================================================
\section{Tiền lương, thù lao và thanh toán}

\subsection{Nguyên tắc chung}

Tiền lương và thù lao được xác định theo:

\begin{itemize}
    \item Hợp đồng hoặc phụ lục hợp đồng.
    \item Loại hình công việc.
    \item Số giờ hoặc số buổi được xác nhận.
    \item Teaching Rate có hiệu lực.
    \item Phụ cấp, thưởng và khoản điều chỉnh hợp lệ.
    \item Thuế, bảo hiểm và nghĩa vụ bắt buộc nếu áp dụng.
\end{itemize}

\subsection{Nhân sự toàn thời gian}

Thu nhập có thể gồm:

\begin{itemize}
    \item Lương cơ bản.
    \item Phụ cấp theo vị trí hoặc điều kiện làm việc.
    \item Tiền làm thêm giờ nếu phát sinh hợp lệ.
    \item Thưởng theo chính sách.
    \item Khoản bổ sung được thỏa thuận.
\end{itemize}

\subsection{Giảng viên và trợ giảng theo buổi}

Giá trị thanh toán có thể được tính theo công thức:

\[
\text{Thù lao giảng dạy}
=
\text{Số đơn vị giảng dạy được xác nhận}
\times
\text{Teaching Rate}
+
\text{Khoản bổ sung hợp lệ}
\]

Teaching Rate áp dụng phải là mức đã được thỏa thuận và có hiệu lực
tại thời điểm buổi học được thực hiện, trừ khi các bên có thỏa thuận khác.

\subsection{Đối soát}

Trước khi thanh toán, các bên có thể đối soát:

\begin{itemize}
    \item Danh sách buổi đã dạy.
    \item Thời lượng thực tế.
    \item Trạng thái hoàn thành.
    \item Buổi bị hủy và bên chịu trách nhiệm.
    \item Công việc chấm bài hoặc hỗ trợ bổ sung.
\end{itemize}

Giảng viên có quyền phản hồi khi phát hiện sai lệch.

% =========================================================
\section{Thưởng và ghi nhận thành tích}

Tiền thưởng không phải khoản mặc nhiên phát sinh,
trừ khi hợp đồng hoặc quy chế có quy định khác.

Tiêu chí thưởng có thể bao gồm:

\begin{itemize}
    \item Chất lượng chuyên môn.
    \item Tỷ lệ hoàn thành công việc.
    \item Mức độ hài lòng hợp lệ của học viên.
    \item Chất lượng phản hồi và hỗ trợ.
    \item Đóng góp phát triển nội dung.
    \item Tuân thủ lịch và quy trình.
\end{itemize}

Không nên quyết định thưởng chỉ dựa trên điểm rating,
vì rating có thể bị ảnh hưởng bởi thiên lệch hoặc hành vi thao túng.

% =========================================================
\section{Nghỉ phép và nghỉ việc riêng}

Nhân sự có thể gửi đơn nghỉ qua hệ thống với các thông tin:

\begin{itemize}
    \item Loại nghỉ.
    \item Thời gian bắt đầu và kết thúc.
    \item Lý do trong phạm vi cần thiết.
    \item Người bàn giao công việc.
    \item Tệp minh chứng nếu trường hợp áp dụng yêu cầu.
\end{itemize}

Quy trình xử lý gồm:

\begin{enumerate}
    \item Gửi yêu cầu.
    \item Kiểm tra số ngày và lịch công việc.
    \item Người có thẩm quyền phê duyệt hoặc từ chối.
    \item Cập nhật chấm công.
    \item Thông báo kết quả cho người yêu cầu.
\end{enumerate}

Nếu vượt quá số ngày nghỉ hưởng lương,
phần vượt có thể được xem xét là nghỉ không hưởng lương
sau khi có thỏa thuận hoặc căn cứ hợp lệ.

Hệ thống không nên tự động chuyển thành nghỉ không lương
nếu chưa hoàn tất bước xác nhận.

% =========================================================
\section{Bàn giao trong thời gian nghỉ}

Giảng viên nghỉ trong thời gian đang phụ trách lớp phải:

\begin{itemize}
    \item Thông báo theo thời hạn hợp lý.
    \item Cập nhật tiến độ giảng dạy.
    \item Bàn giao tài liệu và thông tin cần thiết.
    \item Phối hợp bố trí lịch bù hoặc người thay thế.
    \item Không tự ý hủy lớp mà chưa thông báo.
\end{itemize}

AILMS phải hạn chế ảnh hưởng đến quyền lợi học viên
và thông báo phương án thay thế kịp thời.

% =========================================================
\section{Đánh giá hiệu quả công việc}

\subsection{Nguồn thông tin đánh giá}

Việc đánh giá có thể dựa trên:

\begin{itemize}
    \item Mức độ hoàn thành nhiệm vụ.
    \item Chất lượng nội dung và giảng dạy.
    \item Tiến độ phản hồi và chấm bài.
    \item Tuân thủ lịch làm việc.
    \item Phản hồi của học viên.
    \item Đánh giá của người quản lý.
    \item Khiếu nại đã được xác minh.
    \item Đóng góp cho tổ chức.
\end{itemize}

\subsection{Nguyên tắc đánh giá}

Việc đánh giá phải:

\begin{itemize}
    \item Dựa trên tiêu chí được thông báo.
    \item Xem xét đúng thời kỳ và phạm vi công việc.
    \item Hạn chế ảnh hưởng của đánh giá giả hoặc thiên lệch.
    \item Cho phép nhân sự xem và phản hồi kết quả.
    \item Không chỉ dựa trên quyết định tự động của AI.
\end{itemize}

\subsection{Rating của học viên}

Rating thấp không tự động chứng minh giảng viên vi phạm.

Hệ thống cần xem xét:

\begin{itemize}
    \item Số lượng đánh giá.
    \item Nội dung phản hồi.
    \item Tính xác thực của người đánh giá.
    \item Mức độ và tần suất của vấn đề.
    \item Dữ liệu lịch học và khiếu nại liên quan.
\end{itemize}

% =========================================================
\section{Đào tạo và phát triển chuyên môn}

AILMS có thể tổ chức:

\begin{itemize}
    \item Đào tạo sử dụng hệ thống.
    \item Đào tạo bảo mật và bảo vệ dữ liệu.
    \item Hướng dẫn xây dựng khóa học.
    \item Hướng dẫn đánh giá và phản hồi học viên.
    \item Đào tạo sử dụng AI có trách nhiệm.
    \item Hoạt động cập nhật chuyên môn.
\end{itemize}

Một số chương trình có thể là điều kiện trước khi được cấp
hoặc duy trì quyền giảng dạy.

% =========================================================
\section{Sở hữu trí tuệ}

\subsection{Nguyên tắc xác định quyền}

Quyền sở hữu trí tuệ phải được xác định dựa trên:

\begin{itemize}
    \item Chủ thể tạo ra nội dung.
    \item Nhiệm vụ được giao.
    \item Nguồn lực được sử dụng.
    \item Hợp đồng hoặc thỏa thuận.
    \item Quy định pháp luật áp dụng.
\end{itemize}

Không nên mặc định rằng mọi nội dung giảng viên tạo ra
trong thời gian hợp đồng đều thuộc sở hữu độc quyền của AILMS.

\subsection{Các mô hình quyền}

Hợp đồng có thể quy định:

\begin{itemize}
    \item Giảng viên giữ quyền sở hữu và cấp phép cho AILMS khai thác.
    \item Các bên đồng sở hữu hoặc chia sẻ quyền khai thác.
    \item AILMS sở hữu quyền tài sản đối với nội dung được tạo theo nhiệm vụ.
    \item Quyền được phân chia riêng theo từng loại tài liệu.
\end{itemize}

\subsection{Cam kết của giảng viên}

Giảng viên cam kết:

\begin{itemize}
    \item Có quyền sử dụng tài liệu đã tải lên.
    \item Trích dẫn nguồn khi cần thiết.
    \item Không sử dụng trái phép nội dung của bên thứ ba.
    \item Thông báo nếu tài liệu có giới hạn giấy phép.
    \item Không đưa dữ liệu bí mật của tổ chức khác vào khóa học.
\end{itemize}

\subsection{Sử dụng sau khi hợp đồng chấm dứt}

Quyền tiếp tục khai thác khóa học sau khi hợp đồng chấm dứt
phải căn cứ vào hợp đồng, bao gồm:

\begin{itemize}
    \item Thời hạn tiếp tục khai thác.
    \item Phạm vi chỉnh sửa và phân phối.
    \item Cơ chế chia sẻ doanh thu nếu có.
    \item Quyền sử dụng tên và hình ảnh giảng viên.
    \item Trách nhiệm cập nhật hoặc gỡ bỏ nội dung.
\end{itemize}

% =========================================================
\section{Bảo mật và dữ liệu nội bộ}

Nhân sự không được tiết lộ trái phép:

\begin{itemize}
    \item Hồ sơ và kết quả học tập của học viên.
    \item Hợp đồng, tiền lương và thông tin nhân sự.
    \item Secret key, token và thông tin kết nối.
    \item Mã nguồn hoặc kiến trúc nội bộ.
    \item Kế hoạch kinh doanh và dữ liệu tài chính.
    \item Tài liệu chưa được công bố.
\end{itemize}

Thông tin chỉ được truy cập và sử dụng trong phạm vi công việc.

Nghĩa vụ bảo mật có thể tiếp tục sau khi hợp đồng chấm dứt
theo thời hạn và phạm vi đã thỏa thuận.

% =========================================================
\section{Sử dụng AI trong công việc}

Nhân sự có thể sử dụng AI để hỗ trợ:

\begin{itemize}
    \item Lập bản nháp nội dung.
    \item Tạo câu hỏi hoặc rubric tham khảo.
    \item Tóm tắt tài liệu.
    \item Phân tích dữ liệu trong phạm vi được phép.
    \item Hỗ trợ phản hồi học viên.
\end{itemize}

Nhân sự không được:

\begin{itemize}
    \item Gửi dữ liệu cá nhân hoặc bí mật vào công cụ AI không được phê duyệt.
    \item Công bố nội dung AI mà không kiểm tra.
    \item Dùng AI để tự động quyết định kỷ luật hoặc chấm điểm quan trọng.
    \item Sử dụng AI để giả mạo năng lực hoặc kết quả công việc.
    \item Đưa tài liệu nội bộ vào hệ thống bên ngoài khi chưa được phép.
\end{itemize}

Người sử dụng AI vẫn chịu trách nhiệm đối với kết quả công việc cuối cùng.

% =========================================================
\section{Xung đột lợi ích}

Nhân sự phải khai báo khi có tình huống có thể ảnh hưởng đến tính khách quan, như:

\begin{itemize}
    \item Có quan hệ cá nhân hoặc tài chính với học viên, đối tác hay nhà cung cấp.
    \item Tham gia kinh doanh cạnh tranh trực tiếp.
    \item Đánh giá hoặc phê duyệt công việc của người có quan hệ lợi ích.
    \item Nhận quà hoặc lợi ích có khả năng ảnh hưởng đến quyết định.
\end{itemize}

Việc có xung đột lợi ích không mặc nhiên là vi phạm,
nhưng việc che giấu hoặc lợi dụng xung đột có thể bị xử lý.

% =========================================================
\section{Phòng chống quấy rối và phân biệt đối xử}

AILMS không chấp nhận:

\begin{itemize}
    \item Quấy rối tình dục hoặc quấy rối bằng lời nói, hình ảnh hay hành vi.
    \item Bắt nạt, đe dọa hoặc xúc phạm.
    \item Phân biệt đối xử trái pháp luật.
    \item Trả đũa người phản ánh hoặc người hỗ trợ điều tra.
    \item Lợi dụng quan hệ giảng viên--học viên để gây áp lực.
\end{itemize}

Phản ánh phải được tiếp nhận bảo mật trong phạm vi có thể,
xử lý khách quan và bảo vệ các bên khỏi hành vi trả đũa.

% =========================================================
\section{Vi phạm và xử lý kỷ luật}

\subsection{Nguyên tắc}

Việc xử lý phải:

\begin{itemize}
    \item Có căn cứ và bằng chứng.
    \item Đúng thẩm quyền và trình tự.
    \item Xem xét tính chất, mức độ và hậu quả.
    \item Cho phép người bị phản ánh giải trình.
    \item Không phạt tiền hoặc cắt lương thay cho xử lý kỷ luật trái quy định.
    \item Không áp dụng đồng thời nhiều hình thức trái pháp luật cho cùng một hành vi.
\end{itemize}

\subsection{Biện pháp quản lý}

Tùy tính chất vụ việc, AILMS có thể:

\begin{itemize}
    \item Nhắc nhở hoặc yêu cầu khắc phục.
    \item Yêu cầu đào tạo lại.
    \item Cảnh báo bằng văn bản.
    \item Tạm ngừng phân công lớp.
    \item Tạm khóa một số quyền để bảo vệ dữ liệu.
    \item Thay đổi phân công theo thỏa thuận hợp lệ.
    \item Áp dụng hình thức xử lý theo nội quy và pháp luật.
\end{itemize}

Khóa tài khoản kỹ thuật khẩn cấp nhằm bảo vệ hệ thống
không tự động đồng nghĩa với kết luận kỷ luật.

% =========================================================
\section{Khiếu nại và giải trình}

Nhân sự có quyền phản ánh hoặc khiếu nại về:

\begin{itemize}
    \item Dữ liệu chấm công.
    \item Tiền lương hoặc Teaching Session Payment.
    \item Kết quả đánh giá.
    \item Phân công công việc.
    \item Hành vi quấy rối hoặc phân biệt đối xử.
    \item Quyết định hạn chế quyền hoặc xử lý vi phạm.
    \item Việc sử dụng dữ liệu cá nhân không phù hợp.
\end{itemize}

Yêu cầu nên bao gồm:

\begin{itemize}
    \item Employee Code.
    \item Nội dung và thời điểm phát sinh.
    \item Quyết định hoặc dữ liệu liên quan.
    \item Bằng chứng nếu có.
    \item Phương án đề nghị xử lý.
\end{itemize}

Người xử lý khiếu nại không nên là người có xung đột lợi ích trực tiếp với vụ việc.

% =========================================================
\section{Audit Log và minh bạch nghiệp vụ}

Hệ thống cần ghi Audit Log đối với:

\begin{itemize}
    \item Tạo và sửa hồ sơ nhân sự.
    \item Cấp, thay đổi hoặc thu hồi vai trò.
    \item Thay đổi trạng thái hợp đồng.
    \item Phê duyệt đơn nghỉ.
    \item Điều chỉnh bảng công.
    \item Phê duyệt khoản thanh toán.
    \item Xử lý khiếu nại và vi phạm.
    \item Phê duyệt hoặc đình chỉ khóa học.
\end{itemize}

Audit Log phải ghi nhận hợp lý:

\begin{itemize}
    \item Người thực hiện.
    \item Thời điểm.
    \item Loại hành động.
    \item Đối tượng bị tác động.
    \item Giá trị trước và sau nếu phù hợp.
    \item Lý do hoặc mã yêu cầu liên quan.
\end{itemize}

Không ghi mật khẩu, token đầy đủ hoặc dữ liệu nhạy cảm không cần thiết vào log.

% =========================================================
\section{Bảo vệ hồ sơ nhân sự}

Hồ sơ nhân sự, hợp đồng và dữ liệu thu nhập chỉ được truy cập bởi
người có nhiệm vụ và quyền phù hợp.

AILMS phải áp dụng:

\begin{itemize}
    \item Kiểm soát truy cập theo vai trò.
    \item Ghi nhận lịch sử truy cập và thay đổi quan trọng.
    \item Giới hạn tải xuống và chia sẻ.
    \item Bảo vệ tệp bằng cơ chế xác thực hoặc URL có thời hạn.
    \item Sao lưu và phục hồi phù hợp.
    \item Xóa hoặc lưu trữ theo thời hạn hợp lệ.
\end{itemize}

Nhân sự có thể yêu cầu xem và chỉnh sửa thông tin không chính xác
sau khi hoàn thành bước xác minh cần thiết.

% =========================================================
\section{Chấm dứt hợp đồng và bàn giao}

Khi hợp đồng chấm dứt, nhân sự phải:

\begin{itemize}
    \item Bàn giao công việc, lớp học và tài liệu.
    \item Hoàn trả thiết bị hoặc tài sản được cấp.
    \item Không giữ bản sao dữ liệu nội bộ trái phép.
    \item Hoàn tất đối soát công việc và khoản thanh toán.
    \item Tiếp tục thực hiện nghĩa vụ bảo mật nếu còn hiệu lực.
\end{itemize}

AILMS có trách nhiệm:

\begin{itemize}
    \item Thu hồi quyền truy cập không còn cần thiết.
    \item Bảo toàn dữ liệu học viên và tiến độ lớp.
    \item Bố trí người thay thế khi cần.
    \item Thực hiện nghĩa vụ thanh toán và xác nhận theo quy định.
    \item Không xóa hồ sơ bắt buộc phải lưu giữ.
\end{itemize}

% =========================================================
\section{Sửa đổi chính sách}

AILMS có thể cập nhật chính sách khi:

\begin{itemize}
    \item Pháp luật lao động thay đổi.
    \item Quy trình hoặc cơ cấu tổ chức thay đổi.
    \item Hệ thống bổ sung chức năng mới.
    \item Có yêu cầu tăng cường bảo mật hoặc minh bạch.
\end{itemize}

Thay đổi ảnh hưởng trực tiếp đến quyền lợi và nghĩa vụ
phải được thông báo bằng hình thức phù hợp.

Việc cập nhật chính sách nội bộ không tự động sửa đổi nội dung hợp đồng
nếu pháp luật hoặc thỏa thuận yêu cầu sự đồng ý của các bên.

% =========================================================
\section{Thông tin liên hệ}

Nhân sự có thể gửi yêu cầu tới:

\begin{itemize}
    \item \textbf{Đơn vị phụ trách:} Phòng Nhân sự AILMS.
    \item \textbf{Email nhân sự:}
    \href{mailto:hr@ailms.example}{hr@ailms.example}.
    \item \textbf{Email bảo mật:}
    \href{mailto:security@ailms.example}{security@ailms.example}.
    \item \textbf{Kênh nội bộ:} Chức năng gửi yêu cầu nhân sự trên AILMS.
\end{itemize}

\textcolor{ailmsred}{
\textit{Lưu ý: Thay các địa chỉ thư điện tử mẫu và thông tin tổ chức
bằng dữ liệu chính thức trước khi ban hành.}
}

% =========================================================
\section{Hiệu lực thi hành}

Chính sách có hiệu lực kể từ ngày được phê duyệt và công bố.

Toàn bộ nhân sự, giảng viên, trợ giảng và người quản lý liên quan
có trách nhiệm đọc, hiểu và tuân thủ chính sách.

\vfill

\begin{center}
    \textbf{ĐẠI DIỆN PHÒNG NHÂN SỰ AILMS}\\[1.5cm]
    \textit{Ký và ghi rõ họ tên}
\end{center}

\end{document}